\ExerciseSection

\begin{enumerate}
\def\labelenumi{\arabic{enumi})}
\tightlist
\item
  \begin{enumerate}
  \def\labelenumii{\alph{enumii})}
  \tightlist
  \item
    Let \((x_{n}, y_{n})_{n \in \mathbf{Z}}\) be a family of points of \(\mathbf{R}^{2}\) such that for all \(n \in \mathbf{Z}\) we have \(x_{n} < x_{n+1}\) and
  \end{enumerate}
\end{enumerate}

\[
\frac {y _ {n + 1} - y _ {n}}{x _ {n + 1} - x _ {n}} <   \frac {y _ {n + 2} - y _ {n + 1}}{x _ {n + 2} - x _ {n + 1}}.
\]

Show that, for all integers \(n \in \mathbf{Z}\) , \(m > 0\) , \(p > 0\) , we have

\[
\frac {y _ {n + m} - y _ {n}}{x _ {n + m} - x _ {n}} <   \frac {y _ {n + m + p} - y _ {n}}{x _ {n + m + p} - x _ {n}}.
\]

\begin{enumerate}
\def\labelenumi{\alph{enumi})}
\setcounter{enumi}{1}
\tightlist
\item
  If \(a > 0\) and \(x \neq 0\) , put
\end{enumerate}

\[
f _ {a} (x) = \frac {a ^ {x} - 1}{x}.
\]

Show that, if \(x < y\) , \(x \neq 0\) , \(y \neq 0\) and \(a \neq 1\) , we have \(f_{a}(x) < f_{a}(y)\) {[}prove the relation first for \(x, y\) integral, by means of \(a\) ), then for \(x, y\) rational, and finally in general{]}.

\begin{enumerate}
\def\labelenumi{\alph{enumi})}
\setcounter{enumi}{2}
\item
  Deduce that, for all \(a > 0\) , the function \(f_{a}(x)\) , defined for \(x \neq 0\) , has a limit on the right and a limit on the left as \(x \to 0\) ; show that these two limits have a common value \(\varphi(a)\) , and that \(\varphi(a) \neq 0\) if \(a \neq 1\) . Show that if \(a \neq 1\) , the function \((\log_{a} x) / (x - 1)\) , defined for \(x \neq 1\) , tends to the limit \(1 / \varphi(a)\) as \(x \to 1\) .
\item
  Show that, for all \(a > 0\) and \(b > 0\) , we have \(\varphi(b) = \varphi(a) \log_a b\) (if \(a \neq 1\) ). Deduce that there exists a number \(e\) between 2 and 4 such that \(\varphi(e) = 1\) , and that \(\varphi(a) = \log_e a\) for all \(a > 0\) .
\end{enumerate}

\begin{enumerate}
\def\labelenumi{\arabic{enumi})}
\setcounter{enumi}{1}
\tightlist
\item
  Show, by induction on n, that for all integers n \textgreater{} 0 we have
\end{enumerate}

\[
2 ^ {n} > \frac {1}{2} n (n + 1).
\]

Deduce that, if \(a > 1\) and \(\alpha > 0\) ,

\[
\lim _ {x \rightarrow + \infty} \frac {a ^ {x}}{x ^ {\alpha}} = + \infty , \quad \lim _ {x \rightarrow + \infty} \frac {\log_ {a} x}{x ^ {\alpha}} = 0
\]

(begin by proving the first of these relations for \(a = 2\) and \(\alpha = 1\) ).

